%-----------------------------------
% Kapitel 3: Ist-Analyse des First-Level-Supports im Incident Management
%-----------------------------------
\section{Ist-Analyse des First-Level-Supports im Incident Management}
\label{istAnalyse}

Kapitel~3 dient der Problemidentifikation im Sinne des in Abschnitt~1.3 dargelegten Vorgehens. Da die Arbeit literaturbasiert angelegt ist, betrachtet sie nicht den Prozess eines einzelnen Unternehmens, sondern einen generischen Referenzprozess, der aus den ITIL-Vorgaben zum Incident Management und zum Service Desk zusammengef\"uhrt wird. Diese Zusammenf\"uhrung stellt eine Eigenleistung dar; die einzelnen Prozessbestandteile sind belegt, ihre Verbindung zu einem durchg\"angigen Ablauf ist es nicht. Abschnitt~3.1 beschreibt den Ablauf in drei Phasen, Abschnitt~3.2 \"uberf\"uhrt ihn in ein BPMN-Modell. Abschnitt~3.3 benennt die Kennzahlen, an denen sich der Prozess messen l\"asst, und Abschnitt~3.4 untersucht Schwachstellen sowie daraus folgende Automatisierungspotenziale.

\subsection{Der First-Level-Support-Prozess nach dem ITIL-Referenzmodell}

Der Ablauf l\"asst sich in drei Phasen gliedern. Abschnitt~3.1.1 behandelt den Ticket-Eingang und die manuelle Triage, Abschnitt~3.1.2 die Informationsbeschaffung und L\"osungsfindung und Abschnitt~3.1.3 die Eskalation in den Second-Level-Support.

\subsubsection{Ticket-Eingang und manuelle Triage}

Meldungen erreichen den Service Desk \"uber mehrere Kan\"ale, die sich im Strukturierungsgrad der \"ubermittelten Angaben unterscheiden. Ein Telefonanruf liefert die Schilderung m\"undlich und erlaubt unmittelbare R\"uckfragen, hinterl\"asst jedoch keine verwertbare Datenstruktur. E-Mails und Chat-Nachrichten liegen als Freitext vor, w\"ahrend ein Self-Service-Portal \"uber Formularfelder bereits eine Vorstrukturierung erzwingt.\\footcite[Vgl.][S.~$\\bullet$]{AXELOS2020} Welche Kan\"ale angeboten werden, richtet sich nach der Nutzergruppe und den vereinbarten Servicezeiten.\\footcite[Vgl.][S.~$\\bullet$]{AXELOS2020} Die Zugangswege bestehen nebeneinander, sodass derselbe Sachverhalt je nach gew\"ahltem Kanal in unterschiedlicher Detailtiefe und Form eintrifft.

Die Erfassung \"uberf\"uhrt die Meldung in das ITSM-Tool. Der Agent legt ein Ticket an und dokumentiert Melder, betroffenen Service, Zeitpunkt sowie die Symptombeschreibung, wobei die Angaben der Anwender h\"aufig unvollst\"andig sind oder Vermutungen \"uber die Ursache anstelle beobachtbarer Symptome enthalten.\\footcite[Vgl.][S.~$\\bullet$]{BeimsZiegenbein} Fehlen entscheidungsrelevante Angaben, stellt der Agent R\"uckfragen, was den Vorgang unterbricht und bei asynchronen Kan\"alen Wartezeiten erzeugt. Mit der Anlage des Tickets beginnt zugleich die Messung der vereinbarten Reaktions- und L\"osungszeiten, weshalb der Erfassungszeitpunkt festgehalten wird.\\footcite[Vgl.][S.~$\\bullet$]{AXELOS2020incident}

Anschlie{\ss}end erfolgt die Triage. Die Kategorisierung ordnet das Ticket \"uber einen mehrstufigen Kategorienbaum einem Service und einer St\"orungsklasse zu, woraus sich die zust\"andige Bearbeitungsgruppe ableitet. Die Priorisierung ergibt sich aus der Einsch\"atzung von Auswirkung und Dringlichkeit, die der Agent anhand der geschilderten Umst\"ande vornimmt.\\footcite[Vgl.][S.~$\\bullet$]{BeimsZiegenbein} Beide Schritte beruhen auf dem Erfahrungswissen des einzelnen Bearbeiters, da sich weder Kategorie noch Priorit\"at unmittelbar aus dem Wortlaut der Meldung ergeben. Bei hohem Meldungsaufkommen findet die Triage unter Zeitdruck statt, weil die Annahme weiterer eingehender Kontakte parallel sicherzustellen ist.\\footcite[Vgl.][S.~$\\bullet$]{AXELOS2020incident} Am Ende der Phase liegt ein erfasstes, kategorisiertes und priorisiertes Ticket vor, dessen weitere Bearbeitung der folgende Abschnitt behandelt.

\subsubsection{Informationsbeschaffung und L\"osungsfindung}

Die Erstdiagnose beginnt mit der Analyse der dokumentierten Symptome. Der Agent pr\"uft, ob die Meldung einem bekannten Fehler entspricht, ob weitere Tickets mit gleichem Muster vorliegen und ob ein laufender Major Incident die St\"orung erkl\"art, dessen Bearbeitung eigenen Regeln folgt.\\footcite[Vgl.][S.~$\\bullet$]{AXELOS2020incident} Erg\"anzend ist zu kl\"aren, ob geplante Wartungsarbeiten oder kurz zuvor durchgef\"uhrte \"Anderungen den geschilderten Effekt hervorrufen, da solche Zusammenh\"ange die weitere Diagnose er\"ubrigen.\\footcite[Vgl.][S.~$\\bullet$]{BeimsZiegenbein} Diese Einordnung entscheidet dar\"uber, ob eine eigenst\"andige Bearbeitung erfolgt oder das Ticket einem bereits laufenden Vorgang zugeordnet wird.

F\"ur die L\"osungsfindung stehen mehrere Informationsquellen bereit. Die Wissensdatenbank enth\"alt dokumentierte L\"osungswege, die Known-Error-Datenbank bekannte Fehler samt zugeh\"origer Workarounds, und die Ticket-Historie erlaubt den Abgleich mit \"ahnlich gelagerten F\"allen.\footcites[Vgl.][S.~$\bullet$]{AXELOS2020incident}[Vgl.][S.~$\bullet$]{BeimsZiegenbein} Die Suche erfolgt manuell \"uber Stichworte, wobei der Treffererfolg von der Formulierung der Suchanfrage und von der Pflege der Best\"ande abh\"angt. Bleibt sie ergebnislos, weitet der Agent die Recherche auf Herstellerdokumentation und weitere allgemeine Quellen aus oder zieht erfahrene Kollegen hinzu. Dieser Schritt beansprucht einen erheblichen Teil der Bearbeitungszeit, insbesondere wenn die Symptombeschreibung des Anwenders von der Terminologie der Wissensartikel abweicht.\footcite[Vgl.][S.~$\bullet$]{BeimsZiegenbein}

F\"uhrt die Recherche zu einem Ergebnis, unternimmt der First Level einen L\"osungsversuch. Typisch sind die Anleitung des Anwenders durch einzelne Schritte, die Anwendung dokumentierter Standardl\"osungen sowie Routinevorg\"ange wie das Zur\"ucksetzen eines Kennworts.\footcites[Vgl.][S.~$\bullet$]{AXELOS2020incident}[Vgl.][S.~$\bullet$]{AXELOS2020} F\"ur wiederkehrende Anliegen halten Organisationen Ablaufbeschreibungen vor, die den Agenten anleiten und ein einheitliches Vorgehen unabh\"angig vom einzelnen Bearbeiter sichern. Gelingt die Wiederherstellung, dokumentiert der Agent das Vorgehen im Ticket, informiert den Anwender und schlie{\ss}t den Vorgang nach dessen Best\"atigung ab.\footcite[Vgl.][S.~$\bullet$]{AXELOS2020} Die Dokumentation dient zugleich als Grundlage f\"ur vergleichbare k\"unftige F\"alle, sofern die L\"osung als Wissensartikel \"ubernommen wird. Bleibt der L\"osungsversuch erfolglos, tritt der Prozess in die Eskalation ein.

\subsubsection{Eskalation in den Second-Level-Support}

Die Weitergabe folgt zwei voneinander unabh\"angigen Logiken. Die funktionale Eskalation greift, wenn dem First Level die fachliche Kompetenz oder die technische Berechtigung zur L\"osung fehlt; das Ticket geht dann an eine Fachgruppe des Second- oder Third-Level-Supports.\footcites[Vgl.][S.~$\bullet$]{AXELOS2020incident}[Vgl.][S.~$\bullet$]{BeimsZiegenbein} Die hierarchische Eskalation richtet sich nicht an Spezialisten, sondern an Entscheidungstr\"ager, und erfolgt, wenn eine Verletzung der SLA-Zielzeiten droht, die Auswirkung der St\"orung hoch ist oder zus\"atzliche Ressourcen freigegeben werden m\"ussen.\footcites[Vgl.][S.~$\bullet$]{BeimsZiegenbein}[Vgl.][S.~$\bullet$]{AXELOS2020incident} Beide Formen k\"onnen denselben Vorgang betreffen. Die Kriterien sind im Vorfeld festgelegt, damit die Entscheidung \"uber eine Weitergabe nicht im Einzelfall neu getroffen werden muss.\footcite[Vgl.][S.~$\bullet$]{BeimsZiegenbein}

Der \"Ubergabe liegt die Zuweisung im ITSM-Tool zugrunde, mit der Zust\"andigkeit und Bearbeitungsstand wechseln. Zur \"Ubergabe geh\"ort neben der Zuweisung eine Zusammenfassung des bisherigen Verlaufs, also der gepr\"uften Annahmen, der durchgef\"uhrten Schritte und ihrer Ergebnisse; bei schwerwiegenden St\"orungen tritt eine m\"undliche Abstimmung hinzu.\footcite[Vgl.][S.~$\bullet$]{AXELOS2020incident} Wie aufw\"andig die weitere Bearbeitung ausf\"allt, h\"angt von der Qualit\"at der bis dahin erfassten Informationen ab, denn eine l\"uckenhafte Dokumentation f\"uhrt zu R\"uckfragen beim Anwender und dazu, dass der Second Level Teile der Diagnose erneut durchf\"uhrt.

Die Zust\"andigkeit f\"ur den Vorgang verbleibt beim Service Desk. Er bleibt Eigent\"umer des Tickets, verfolgt den Bearbeitungsstand, informiert den Anwender \"uber Zwischenst\"ande und nimmt nach der L\"osung dessen Best\"atigung entgegen.\footcite[Vgl.][S.~$\bullet$]{AXELOS2020} Der so beschriebene Ablauf bildet den Referenzprozess dieser Arbeit. Abschnitt~3.2 \"uberf\"uhrt ihn in ein BPMN-Modell.

\subsection{Modellierung des Ist-Zustandes (BPMN-Modell)}

% Hier Inhalt einfuegen

\subsection{Kennzahlen zur Prozessbewertung}

% Hier Inhalt einfuegen

\subsection{Schwachstellenanalyse und Identifikation von Automatisierungspotenzialen}

\subsubsection{Zeitverluste durch manuelle Kategorisierung}

% Hier Inhalt einfügen

\subsubsection{Ineffizienzen durch unstrukturierte Nutzereingaben}

% Hier Inhalt einfügen
